% TODO checklist:
% - [x] README.md: Platform goals, feature list
% - [x] Q&A.md: Q1-Q5 (Platform capabilities), evaluation criteria discussions
% - [x] sapthesis-doc.tex: Abstract and contributions framing

\section{Research Objectives}
\label{sec:objectives}

This thesis pursues a primary objective supported by several sub-objectives that collectively address the challenges identified in Section~\ref{sec:problem-statement}. The objectives are formulated to be specific, measurable, and verifiable through the evaluation presented in Chapter~\ref{ch:evaluation}.

\subsection{Primary Objective}

\begin{quote}
\textbf{Primary Objective:} \emph{Design and implement an enterprise-grade AI agent orchestration platform that integrates comprehensive security, multi-tenancy, graph database querying, observability, and extensible tool capabilities into a production-ready system suitable for deployment at research institutions and enterprises.}
\end{quote}

This objective encompasses the full lifecycle from architectural design through implementation to production deployment, with the Cineca Agentic Platform serving as the concrete realization.

\subsection{Sub-Objectives}

The primary objective decomposes into six sub-objectives, each addressing a specific challenge area.

\subsubsection{SO1: Security and Governance Architecture}

\paragraph{Objective:} Design and implement a comprehensive security framework providing:

\begin{enumerate}[label=\alph*)]
    \item OIDC/JWT-based authentication with support for federated identity providers
    \item Role-Based Access Control (RBAC) with scope-based permissions
    \item Per-tool authorization policies integrated with the RBAC system
    \item Comprehensive audit logging for compliance and forensics
    \item PII detection and redaction capabilities
    \item Intent-based filtering to block potentially dangerous operations
\end{enumerate}

\paragraph{Success Criteria:}
\begin{itemize}
    \item All API endpoints enforce authentication and authorization
    \item Audit logs capture all significant operations with sufficient context
    \item PII scrubbing correctly identifies and redacts sensitive data patterns
    \item Security mechanisms do not introduce unacceptable latency overhead ($<$50ms per request)
\end{itemize}

\subsubsection{SO2: Multi-Tenancy Implementation}

\paragraph{Objective:} Implement native multi-tenancy with:

\begin{enumerate}[label=\alph*)]
    \item Database-level tenant isolation for all data entities
    \item Per-tenant configuration defaults (models, rate limits, policies)
    \item Tenant-scoped API endpoints with header-based tenant identification
    \item Administrative delegation within tenant boundaries
    \item Cross-tenant data leakage prevention at all system layers
\end{enumerate}

\paragraph{Success Criteria:}
\begin{itemize}
    \item Tenant A cannot access Tenant B's data through any API path
    \item Per-tenant configurations correctly override system defaults
    \item Tenant isolation is enforced at database query level (not just application level)
    \item Multi-tenant operations scale linearly with tenant count
\end{itemize}

\subsubsection{SO3: Graph Database Integration with NL-to-Cypher}

\paragraph{Objective:} Develop a Natural Language to Cypher pipeline enabling:

\begin{enumerate}[label=\alph*)]
    \item Translation of natural language questions to syntactically correct Cypher queries
    \item Multi-layer safety validation (syntax, read-only, tenant boundaries, resource limits)
    \item Query execution against Memgraph graph database
    \item Result summarization into natural language responses
    \item Deterministic test mode for reliable pipeline testing
\end{enumerate}

\paragraph{Success Criteria:}
\begin{itemize}
    \item Generated Cypher queries pass syntax validation at $>$95\% rate
    \item Safety validation correctly rejects all mutation attempts
    \item Tenant filters are automatically injected into all generated queries
    \item Query execution respects configured timeout and result limits
    \item Test mode enables deterministic, LLM-independent testing
\end{itemize}

\subsubsection{SO4: Operational Observability Stack}

\paragraph{Objective:} Implement comprehensive observability through:

\begin{enumerate}[label=\alph*)]
    \item Prometheus-compatible metrics covering HTTP requests, agent runs, jobs, tools, and LLM calls
    \item Structured logging with request ID correlation across all components
    \item OpenTelemetry distributed tracing for end-to-end request visualization
    \item Kubernetes-compatible health probes (liveness, readiness, startup)
    \item Component-level health checks with degradation detection
\end{enumerate}

\paragraph{Success Criteria:}
\begin{itemize}
    \item All key operations are instrumented with timing and outcome metrics
    \item Log entries include correlation IDs enabling request tracing
    \item Trace spans correctly propagate across async boundaries
    \item Health probes accurately reflect system readiness
    \item Degraded states (e.g., provider unavailability) are correctly reported
\end{itemize}

\subsubsection{SO5: LLM Provider Resilience and Management}

\paragraph{Objective:} Design an LLM provider management system featuring:

\begin{enumerate}[label=\alph*)]
    \item Provider abstraction supporting multiple LLM backends (OpenAI, Ollama, Azure)
    \item Circuit breaker pattern for provider failure isolation
    \item Automatic provider fallback with configurable chains
    \item Cost tracking and budget management
    \item Default Model Resolver with hierarchy (user → tenant → system)
    \item Provider health monitoring and reporting
\end{enumerate}

\paragraph{Success Criteria:}
\begin{itemize}
    \item Provider failures are isolated within circuit breaker timeouts
    \item Fallback chains automatically route to healthy providers
    \item Cost estimates are computed for all LLM calls
    \item Model defaults correctly resolve through the configuration hierarchy
    \item Provider health status is exposed via metrics and health endpoints
\end{itemize}

\subsubsection{SO6: Extensible Tool Ecosystem}

\paragraph{Objective:} Develop an MCP-compatible tool framework supporting:

\begin{enumerate}[label=\alph*)]
    \item Tool registration with schema-based input/output validation
    \item Tool discovery API with filtering by category and authorization
    \item Per-tool RBAC with scope requirements
    \item Comprehensive tool invocation logging
    \item Straightforward patterns for adding new tools
    \item Test fixtures for tool development
\end{enumerate}

\paragraph{Success Criteria:}
\begin{itemize}
    \item Tool discovery returns only tools the caller is authorized to invoke
    \item Input validation rejects malformed invocation requests
    \item Tool invocations are fully logged with inputs, outputs, and timing
    \item New tools can be added by following documented patterns
    \item Tool test coverage exceeds 80\% across all tool categories
\end{itemize}

\subsection{Non-Objectives}

To clarify the scope of this work, the following are explicitly \emph{not} objectives:

\begin{itemize}
    \item Advancing the state of the art in LLM architectures or training methods
    \item Developing novel natural language understanding algorithms
    \item Creating a general-purpose workflow orchestration engine
    \item Building a comprehensive RAG (Retrieval-Augmented Generation) library
    \item Replacing existing agent frameworks for all use cases
\end{itemize}

The focus remains on the integration, orchestration, and operationalization of existing AI capabilities within an enterprise-grade platform architecture.

\subsection{Objective Traceability}

Table~\ref{tab:objective-traceability} maps each sub-objective to the challenges it addresses and the evaluation sections where it is verified.

\begin{table}[ht]
\centering
\caption{Objective traceability matrix.}
\label{tab:objective-traceability}
\small
\begin{tabular}{llll}
\hline
\textbf{Sub-Objective} & \textbf{Addresses Challenge} & \textbf{Primary Chapter} & \textbf{Evaluation} \\
\hline
SO1: Security & Challenge 1 & Chapter~\ref{ch:security} & Sec.~\ref{sec:security-eval} \\
SO2: Multi-Tenancy & Challenge 2 & Chapter~\ref{ch:security} & Sec.~\ref{sec:functional-eval} \\
SO3: NL-to-Cypher & Challenge 3 & Chapter~\ref{ch:orchestration} & Sec.~\ref{sec:functional-eval} \\
SO4: Observability & Challenge 4 & Chapter~\ref{ch:observability} & Sec.~\ref{sec:performance-eval} \\
SO5: LLM Resilience & Challenge 5 & Chapter~\ref{ch:orchestration} & Sec.~\ref{sec:availability-ft} \\
SO6: Tool Ecosystem & Challenge 6 & Chapter~\ref{ch:mcp-tools} & Sec.~\ref{sec:functional-eval} \\
\hline
\end{tabular}
\end{table}

\subsection{Validation Approach}

Each objective will be validated through a combination of:

\begin{description}
    \item[Functional Testing:] Automated tests verifying that implemented features meet specifications. The test suite includes over 2,700 test functions covering all major subsystems.
    
    \item[Performance Benchmarking:] Quantitative measurements of latency, throughput, and resource utilization under realistic workloads.
    
    \item[Security Assessment:] Evaluation of authentication, authorization, and data protection mechanisms against defined threat scenarios.
    
    \item[Operational Evaluation:] Assessment of monitoring, alerting, and operational procedures in production-like environments.
    
    \item[Comparative Analysis:] Comparison with existing frameworks to demonstrate relative capabilities and limitations.
\end{description}

The evaluation chapter (Chapter~\ref{ch:evaluation}) presents detailed results for each validation approach.

